% Abstract.  One paragraph, one claim: a sampling rate of order log m / m suffices for
% the sub-sampled Fiedler split to reproduce the full-data split, and the imbalance eta
% is the parameter that governs the cost.
\begin{abstract}
Top-down spectral reconstruction of a latent binary phylogeny spends $\Theta(m^2)$ on pairwise similarities at its first split, before the recursion has taken a step. Sampling a Bernoulli subset of the entries removes that cost, but phylogenies are imbalanced, and the guarantees available for recursive spectral bipartitioning are stated for clusters of comparable size.
We treat recovery from a sample as an eigenvector perturbation problem. The split of the full similarity matrix is recovered exactly from the sign pattern of the Fiedler vector of a sparse sample, and that split is an edge cut of the tree for any topology.

A sampling rate $p=\Theta(\log m/m)$ suffices at fixed imbalance. Working coordinate-wise, through a Neumann expansion and a per-node Bernstein bound rather than a Davis-Kahan $\ell_2$ transfer, removes a $\sqrt m$ factor from the entry-wise error bound.
The imbalance ratio $\eta$ governs the result and enters twice: it raises the coherence of the leading subspace, which raises the sampling rate recovery needs, and it shrinks the spectral gap, which can close recovery off entirely. The rate depends on $\eta$ as $(1+\eta)^3$.
Simulations on closed-form block models and on Kingman trees with JC69 sequences reproduce the predicted $\Theta(\log m/m)$ scaling and separate the two mechanisms.
\end{abstract}
